\textbf{Proof.}
Put \(m=m_1\wedge m_0\) and let \(W_1\) denote one-dimensional Wasserstein distance.  For laws \(R,S\) on \([0,1]\), layer cake and Tonelli give
\[
W_1(R,S)=\int_0^1|R^{-1}(u)-S^{-1}(u)|\,du
=\int_0^1|R(y)-S(y)|\,dy.
\tag{1}
\]
For each arm, Jensen's inequality and the binomial variance identity imply
\[
\mathbb E_PW_1(\widehat F_a,F_a)
\le\int_0^1\sqrt{\frac{F_a(y)\{1-F_a(y)\}}{m_a}}\,dy
\le\frac1{2\sqrt{m_a}}.
\tag{2}
\]
Because \(y\) and \(y^2\) are one- and two-Lipschitz on \([0,1]\),
\[
|\widehat\mu_a-\mu_a|\le W_1(\widehat F_a,F_a),
\qquad
|\widehat\sigma_a^2-\sigma_a^2|\le4W_1(\widehat F_a,F_a).
\tag{3}
\]
The inequality \(|xy-x'y'|\le|x-x'|+|y-y'|\), applied to the endpoint quantile products, gives
\[
|\widehat c_e-c_e|
\le2\{W_1(\widehat F_1,F_1)+W_1(\widehat F_0,F_0)\},
\qquad e\in\{L,U\}.
\tag{4}
\]
Every block of \(\widehat Q_e-Q(c_e)\) is a linear combination of \(I_n\) and \(\mathbf1\mathbf1^{\top}\), with coefficients controlled by (3)--(4).  Since their operator norms are \(1\) and \(n\), respectively, (2)--(4) yield a numerical constant \(K_M<\infty\) such that
\[
\mathbb E_PM_n\le K_M\frac n{\sqrt m}.
\tag{5}
\]

For \(B\in\mathcal B_{\kappa,\overline a}\), let \(Z_B=\Omega^{1/2}B\).  Regularity makes \(Z_B^{\top}Z_B\) invertible, and
\[
P_B=Z_B(Z_B^{\top}Z_B)^{-1}Z_B^{\top}
\tag{6}
\]
is an orthogonal projector.  Hence \(0\preceq P_B\preceq I_{2n}\), and congruence gives
\[
0\preceq A_B
=\Pi^{-1}\Omega^{1/2}P_B\Omega^{1/2}\Pi^{-1}
\preceq A_{\mathrm{HT}}=\Pi^{-1}\Omega\Pi^{-1}.
\tag{7}
\]
Both matrices are positive semidefinite.  Therefore their nuclear norms are their traces, and
\[
\lVert A_B\rVert_*
\le\operatorname{tr}(A_{\mathrm{HT}})
=\sum_{i,a}\frac{1-\pi_{i,a}}{\pi_{i,a}}
\le2\underline\pi^{-1}n^{1+\beta}.
\tag{8}
\]
The same last bound holds for \(A_{\mathrm{HT}}\).  Substitution in the definition of \(C_n\) proves
\[
C_n\le4\underline\pi^{-1}n^{\beta-1}.
\tag{9}
\]
Trace duality and the one-Lipschitz property of the endpoint maximum prove
\[
\sup_{B\in\mathcal B_{\kappa,\overline a}}|\widehat w(B)-w(B)|\le C_nM_n.
\tag{10}
\]

For the zero-covariance matrix, write \(\Delta_e=\widehat Q_e-Q(c_e)\) and \(d_e=\widehat c_e-c_e\).  Affinity gives
\[
\widehat Q(0)-Q(0)=\Delta_e+d_eE.
\tag{11}
\]
Choose a unit vector \(z\perp\mathbf1\), possible because \(n\ge3\).  Evaluating the off-diagonal block of \(\Delta_e\) between \((z,0)\) and \((0,z)\) shows
\[
|d_e|\le\lVert\Delta_e\rVert_{\mathrm{op}}.
\tag{12}
\]
Since \(\lVert E\rVert_{\mathrm{op}}=1\), (11)--(12) yield
\[
\lVert\widehat Q(0)-Q(0)\rVert_{\mathrm{op}}\le2M_n.
\tag{13}
\]
Trace duality and (8) now give
\[
\sup_{B\in\mathcal B_{\kappa,\overline a}}|\widehat g_B(0)-g_B(0)|
\le2n^{-2}\sup_B\lVert A_B\rVert_*M_n
\le2C_nM_n.
\tag{14}
\]
The exact premise \(nC_n\mathbb E_PM_n\to0\) makes the expectations of (10) and (14) \(o(n^{-1})\).  Finally, (5) and (9) show
\[
nC_n\mathbb E_PM_n
\le4\underline\pi^{-1}K_M\frac{n^{1+\beta}}{\sqrt{m_1\wedge m_0}},
\]
proving the stated sufficient pilot-size condition.  No known-exposure or dependency-neighborhood property was used. \(\square\)